% Gesamttest «Terme und Faktorisieren» — Quelle des PDFs (python3 scripts/build-lp-pdf.py terme-faktorisieren).
% Musterlösung und Raster: bewertungspaket.tex. Alle Werte mit python3 nachgerechnet
% (scripts/lp/terme-faktorisieren/zahlen.py). Nach der Prüfung vom 10.10.2026: G1 verschachtelt mit Probe, G3(a) Flächenmodell,
% G5 Quadrat-Skizze (jedes Kapitelziel geprüft), G8 neu (nicht mehr das Muster k(x² − c²) aus Festhalten und 4e). 25 Punkte, rund 30 Minuten. Kein \lt/\gt — pdfLaTeX kennt sie nicht.
\documentclass[11pt]{article}
\usepackage{../lp-druck}
\renewcommand{\lpname}{Terme und Faktorisieren}
\renewcommand{\lpteilgebiet}{Arithmetik/Algebra ·}
\renewcommand{\lpfuss}{Gesamttest · 25 Punkte}
\begin{document}
\lpkopf{Gesamttest}{%
  \vspace{4pt}\noindent\begin{tabularx}{\linewidth}{@{}X X@{}}
  Code (kein Name): \hrulefill & Punkte: \hrulefill\ / 25
  \end{tabularx}}

\begin{lpinfo}{Anleitung}
\textbf{Zeit:} rund 30 Minuten. \textbf{Hilfsmittel:} keine, kein Taschenrechner.
Schreib zu jeder Aufgabe den \textbf{Rechenweg} auf; bewertet wird der Weg. Gib Ergebnisse \textbf{vollständig vereinfacht}
an, beim Faktorisieren \textbf{vollständig faktorisiert} (als Produkt).
Danach: selbst mit dem \emph{Bewertungspaket} bewerten oder die Lösung samt diesem Blatt scannen bzw.\ fotografieren
(ohne Namen und Standort) und mit dem Bewertungspaket einer KI geben.
\end{lpinfo}

\lpteil{Teil A · Terme vereinfachen}{Kapitel 1 · 5 Punkte}

\begin{lpaufgabe}{G1}{Verschachtelte Klammern}{3 P}
Vereinfache: $\;5x - 2\bigl(3x - (x - 4)\bigr)$. Mach danach die Probe mit $x = 2$.
\lpplatz{6}
\end{lpaufgabe}

\begin{lpaufgabe}{G2}{Minus vor der Klammer}{2 P}
Vereinfache: $\;3a^2 - (a^2 - 2a) - 2a$
\lpplatz{3}
\end{lpaufgabe}

\lpteil{Teil B · Ausmultiplizieren und binomische Formeln}{Kapitel 2 und 3 · 8 Punkte}

\begin{lpaufgabe}{G3}{Klammer mal Klammer}{3 P}
\begin{enumerate}[label=(\alph*),leftmargin=*,itemsep=3pt]
\item Zeichne das Flächenmodell zu $(x + 5)(x + 1)$ und schreib das Produkt als Summe.
\item Multipliziere aus und fasse zusammen: $\;(2x + 5)(x - 3)$
\end{enumerate}
\lpplatz{7}
\end{lpaufgabe}

\begin{lpaufgabe}{G4}{Binomische Formel}{3 P}
Rechne mit einer binomischen Formel: $\;(3y - 4)^2$
\lpplatz{4}
\end{lpaufgabe}

\begin{lpaufgabe}{G5}{Kopfrechnen am Quadrat}{2 P}
Berechne $\;31^2\;$ mit einer binomischen Formel, ohne schriftliches Multiplizieren. Zeichne dazu ein Quadrat mit der
Seite $30 + 1$ und zeig darin, woher das mittlere Glied kommt.
\lpplatz{6}
\end{lpaufgabe}

\lpteil{Teil C · Faktorisieren}{Kapitel 4 · 12 Punkte}

\begin{lpaufgabe}{G6}{Ausklammern}{3 P}
Klammere so viel wie möglich aus: $\;12a^2 + 18a$
\lpplatz{3}
\end{lpaufgabe}

\begin{lpaufgabe}{G7}{Zweiklammersatz}{3 P}
Faktorisiere: $\;x^2 - 2x - 24$
\lpplatz{4}
\end{lpaufgabe}

\begin{lpaufgabe}{G8}{Vollständig faktorisieren}{3 P}
Faktorisiere vollständig: $\;3x^2 + 6x + 3$
\lpplatz{4}
\end{lpaufgabe}

\begin{lpaufgabe}{G9}{Fehler finden}{3 P}
Jemand soll $\;x^2 + 6x + 9\;$ faktorisieren und schreibt:
\[ x^2 + 6x + 9 = x(x + 6) + 9 \]
Warum ist das nicht faktorisiert? Begründe in einem Satz. Faktorisiere dann richtig und mach die Probe durch Ausmultiplizieren.
\lpplatz{6}
\end{lpaufgabe}
\end{document}
